\documentclass[11pt,a4paper]{article}

\usepackage[utf8]{inputenc}
\usepackage[T1]{fontenc}
\usepackage[french]{babel}
\usepackage[a4paper,margin=2.3cm]{geometry}
\usepackage{amsmath,amssymb,amsthm,mathtools}
\usepackage{lmodern}
\usepackage{enumitem}
\usepackage{hyperref}
\usepackage{xcolor}
\usepackage{fancyhdr}
\usepackage{stmaryrd}
\usepackage{mathrsfs}

\hypersetup{
    colorlinks=true,
    linkcolor=black,
    urlcolor=blue,
    citecolor=black
}

\setlength{\parindent}{0pt}
\setlength{\parskip}{0.6em}

\pagestyle{fancy}
\fancyhf{}
\fancyhead[L]{Algèbre linéaire --- Chapitre 1 PGB}
\fancyhead[R]{Vocabulaire ensembliste et raisonnement}
\fancyfoot[C]{\thepage}

% Environnements
\newtheorem{definition}{Définition}[section]
\newtheorem{theoreme}[definition]{Théorème}
\newtheorem{proposition}[definition]{Proposition}
\newtheorem{corollaire}[definition]{Corollaire}
\newtheorem{lemme}[definition]{Lemme}
\theoremstyle{remark}
\newtheorem{remarque}[definition]{Remarque}
\newtheorem{exemple}[definition]{Exemple}
\newtheorem*{methode}{Méthode}
\newtheorem*{idee}{Idée}

% Commandes
\newcommand{\R}{\mathbb{R}}
\newcommand{\N}{\mathbb{N}}
\newcommand{\Z}{\mathbb{Z}}
\newcommand{\Q}{\mathbb{Q}}
\newcommand{\C}{\mathbb{C}}
\newcommand{\PP}{\mathcal{P}}

\begin{document}

\begin{center}
    {\LARGE \textbf{Chapitre 1 --- Vocabulaire ensembliste et méthodes de raisonnement}}\\[0.5em]
    {\large Cours fondamental pour l’algèbre linéaire}
\end{center}

\hrule
\vspace{1em}

\tableofcontents

\newpage

\section{Pourquoi ce chapitre est fondamental ?}

Avant de manipuler des espaces vectoriels, des applications linéaires, des matrices ou des déterminants, il faut maîtriser deux choses :
\begin{itemize}
    \item le \textbf{langage des ensembles}, car une grande partie des objets mathématiques sont des ensembles munis de propriétés ;
    \item les \textbf{méthodes de raisonnement}, car une preuve correcte ne consiste pas à faire des calculs au hasard, mais à enchaîner de la logique.
\end{itemize}

Dans tout le cours d’algèbre linéaire, on rencontrera des formulations comme :
\begin{itemize}
    \item « montrer que \(F\) est un sous-espace vectoriel de \(E\) » ;
    \item « montrer que deux ensembles sont égaux » ;
    \item « montrer qu’une décomposition existe et est unique » ;
    \item « prouver qu’une famille est libre » ;
    \item « raisonner par récurrence sur la dimension ».
\end{itemize}

Ce chapitre constitue donc la base logique de tout le reste.

\section{Notion d’ensemble}

\subsection{Définition intuitive}

\begin{definition}
Un \textbf{ensemble} est une collection d’objets, appelés \textbf{éléments} de cet ensemble.
\end{definition}

On note généralement les ensembles par des lettres majuscules \(A,B,E,F,\dots\), et les éléments par des lettres minuscules \(x,y,a,b,\dots\).

Si \(x\) est un élément de \(A\), on écrit
\[
x\in A.
\]
Si \(x\) n’est pas un élément de \(A\), on écrit
\[
x\notin A.
\]

\begin{exemple}
\begin{itemize}
    \item \(2\in \N\),
    \item \(-3\in \Z\),
    \item \sqrt{2}\in \R\),
    \item \(i\in \C\),
    \item \(\pi \notin \Q\).
\end{itemize}
\end{exemple}

\subsection{Ensembles usuels}

On utilisera constamment :
\[
\N=\{0,1,2,3,\dots\},
\quad
\Z=\{\dots,-2,-1,0,1,2,\dots\},
\]
\[
\Q=\left\{\frac{p}{q}\,;\,p\in\Z,\ q\in\Z\setminus\{0\}\right\},
\quad
\R=\text{ensemble des nombres réels},
\quad
\C=\text{ensemble des nombres complexes}.
\]

\begin{remarque}
Suivant les conventions, certains auteurs définissent \(\N=\{1,2,3,\dots\}\). Dans ce cours, on adopte la convention usuelle en analyse et en algèbre linéaire :
\[
0\in \N.
\]
\end{remarque}

\subsection{Définition par compréhension et par extension}

\begin{definition}
Un ensemble peut être défini :
\begin{itemize}
    \item \textbf{par extension}, en listant ses éléments ;
    \item \textbf{par compréhension}, en donnant une propriété caractéristique.
\end{itemize}
\end{definition}

\begin{exemple}
\[
A=\{1,2,3,4\}
\]
est une définition par extension.

L’ensemble des entiers pairs peut s’écrire
\[
E=\{n\in\Z \mid \exists k\in\Z,\ n=2k\}.
\]
C’est une définition par compréhension.
\end{exemple}

\begin{remarque}
Dans l’écriture
\[
\{x\in E \mid P(x)\},
\]
le symbole \(P(x)\) désigne une propriété dépendant de \(x\). L’ensemble ainsi défini est l’ensemble des éléments de \(E\) qui vérifient \(P(x)\).
\end{remarque}

\section{Inclusion et égalité d’ensembles}

\subsection{Inclusion}

\begin{definition}
Soient \(A\) et \(B\) deux ensembles. On dit que \(A\) est inclus dans \(B\), et on note
\[
A\subset B,
\]
si tout élément de \(A\) appartient à \(B\), c’est-à-dire
\[
\forall x,\ x\in A \Longrightarrow x\in B.
\]
\end{definition}

\begin{remarque}
L’inclusion \(A\subset B\) ne signifie pas forcément que \(A\) est strictement plus petit que \(B\). On peut avoir \(A=B\) et donc \(A\subset B\).
\end{remarque}

\begin{exemple}
\[
\N\subset \Z\subset \Q\subset \R\subset \C.
\]
\end{exemple}

\subsection{Égalité de deux ensembles}

\begin{proposition}
Deux ensembles \(A\) et \(B\) sont égaux si et seulement si
\[
A\subset B
\quad \text{et} \quad
B\subset A.
\]
\end{proposition}

\begin{proof}
Supposons d’abord que \(A=B\). Alors tout élément de \(A\) est un élément de \(B\), donc \(A\subset B\). De même, \(B\subset A\).

Réciproquement, supposons que \(A\subset B\) et \(B\subset A\). Alors tout élément de \(A\) appartient à \(B\), et tout élément de \(B\) appartient à \(A\). Les deux ensembles ont donc exactement les mêmes éléments, ce qui signifie que
\[
A=B.
\]
\end{proof}

\begin{methode}
Pour montrer une égalité d’ensembles
\[
A=B,
\]
on procède presque toujours en deux étapes :
\[
A\subset B
\quad \text{puis} \quad
B\subset A.
\]
C’est une méthode essentielle.
\end{methode}

\begin{exemple}
Montrons que
\[
\{x\in\R \mid x^2=1\}=\{-1,1\}.
\]

\textbf{Première inclusion :} soit \(x\in\R\) tel que \(x^2=1\). Alors
\[
x^2-1=0 \Longleftrightarrow (x-1)(x+1)=0,
\]
donc \(x=1\) ou \(x=-1\). Ainsi \(x\in\{-1,1\}\), donc
\[
\{x\in\R \mid x^2=1\}\subset \{-1,1\}.
\]

\textbf{Seconde inclusion :} si \(x\in\{-1,1\}\), alors \(x=1\) ou \(x=-1\), et dans les deux cas \(x^2=1\). Donc
\[
\{-1,1\}\subset \{x\in\R \mid x^2=1\}.
\]

On conclut que
\[
\{x\in\R \mid x^2=1\}=\{-1,1\}.
\]
\end{exemple}

\section{Opérations sur les ensembles}

\subsection{Union, intersection, différence}

\begin{definition}
Soient \(A\) et \(B\) deux ensembles.
\begin{itemize}
    \item L’\textbf{union} de \(A\) et \(B\) est l’ensemble
    \[
    A\cup B=\{x \mid x\in A \text{ ou } x\in B\}.
    \]
    \item L’\textbf{intersection} de \(A\) et \(B\) est l’ensemble
    \[
    A\cap B=\{x \mid x\in A \text{ et } x\in B\}.
    \]
    \item La \textbf{différence} de \(A\) et \(B\) est l’ensemble
    \[
    A\setminus B=\{x\in A \mid x\notin B\}.
    \]
\end{itemize}
\end{definition}

\begin{exemple}
Si
\[
A=\{1,2,3\},
\qquad
B=\{3,4,5\},
\]
alors
\[
A\cup B=\{1,2,3,4,5\},
\qquad
A\cap B=\{3\},
\qquad
A\setminus B=\{1,2\}.
\]
\end{exemple}

\subsection{Complémentaire}

\begin{definition}
Soit \(E\) un ensemble de référence, et \(A\subset E\). Le \textbf{complémentaire} de \(A\) dans \(E\), noté \(E\setminus A\) ou parfois \(A^c\), est l’ensemble des éléments de \(E\) qui n’appartiennent pas à \(A\) :
\[
E\setminus A=\{x\in E\mid x\notin A\}.
\]
\end{definition}

\begin{exemple}
Dans \(E=\R\), le complémentaire de \([0,1]\) est
\[
\R\setminus [0,1]=]-\infty,0[\ \cup\ ]1,+\infty[.
\]
\end{exemple}

\subsection{Disjonction}

\begin{definition}
Deux ensembles \(A\) et \(B\) sont dits \textbf{disjoints} si
\[
A\cap B=\varnothing,
\]
où \(\varnothing\) désigne l’ensemble vide.
\end{definition}

\begin{exemple}
\[
]-\infty,0[ \quad \text{et} \quad [0,+\infty[
\]
sont pas disjoints car \(0\in [0,+\infty[\), mais \(0\notin ]-\infty,0[\).

En revanche,
\[
]-\infty,0]
\quad \text{et} \quad
[0,+\infty[
\]
ne sont pas disjoints.
\end{exemple}

\section{Propriétés fondamentales des opérations ensemblistes}

\begin{proposition}[Commutativité]
Pour tous ensembles \(A,B\),
\[
A\cup B=B\cup A
\quad \text{et} \quad
A\cap B=B\cap A.
\]
\end{proposition}

\begin{proof}
Montrons par exemple que \(A\cup B=B\cup A\).

Soit \(x\in A\cup B\). Alors \(x\in A\) ou \(x\in B\). Mais l’énoncé « \(x\in A\) ou \(x\in B\) » est équivalent à « \(x\in B\) ou \(x\in A\) ». Donc \(x\in B\cup A\). Ainsi
\[
A\cup B\subset B\cup A.
\]
De même, on montre
\[
B\cup A\subset A\cup B.
\]
Donc
\[
A\cup B=B\cup A.
\]

La preuve pour l’intersection est analogue.
\end{proof}

\begin{proposition}[Associativité]
Pour tous ensembles \(A,B,C\),
\[
(A\cup B)\cup C=A\cup(B\cup C),
\]
\[
(A\cap B)\cap C=A\cap(B\cap C).
\]
\end{proposition}

\begin{proof}
Montrons la première égalité.

Soit \(x\). On a
\[
x\in (A\cup B)\cup C
\Longleftrightarrow (x\in A\cup B)\ \text{ou}\ x\in C
\]
\[
\Longleftrightarrow (x\in A \text{ ou } x\in B)\ \text{ou}\ x\in C
\]
\[
\Longleftrightarrow x\in A \text{ ou } (x\in B \text{ ou } x\in C)
\]
\[
\Longleftrightarrow x\in A \text{ ou } x\in B\cup C
\Longleftrightarrow x\in A\cup(B\cup C).
\]
Les deux ensembles ont donc les mêmes éléments.
\end{proof}

\begin{proposition}[Distributivité]
Pour tous ensembles \(A,B,C\),
\[
A\cap(B\cup C)=(A\cap B)\cup(A\cap C),
\]
\[
A\cup(B\cap C)=(A\cup B)\cap(A\cup C).
\]
\end{proposition}

\begin{proof}
Montrons la première formule.

Soit \(x\). Alors
\[
x\in A\cap(B\cup C)
\Longleftrightarrow x\in A \text{ et } (x\in B \text{ ou } x\in C).
\]
Or
\[
P\text{ et }(Q\text{ ou }R)
\Longleftrightarrow
(P\text{ et }Q)\text{ ou }(P\text{ et }R).
\]
Donc
\[
x\in A\cap(B\cup C)
\Longleftrightarrow
(x\in A \text{ et } x\in B)\text{ ou }(x\in A \text{ et } x\in C)
\]
\[
\Longleftrightarrow x\in (A\cap B)\cup(A\cap C).
\]
Donc
\[
A\cap(B\cup C)=(A\cap B)\cup(A\cap C).
\]
\end{proof}

\begin{proposition}[Formules de De Morgan]
Soit \(E\) un ensemble de référence, et \(A,B\subset E\). Alors
\[
E\setminus(A\cup B)=(E\setminus A)\cap(E\setminus B),
\]
\[
E\setminus(A\cap B)=(E\setminus A)\cup(E\setminus B).
\]
\end{proposition}

\begin{proof}
Montrons la première égalité.

Soit \(x\in E\). Alors
\[
x\in E\setminus(A\cup B)
\Longleftrightarrow x\notin A\cup B
\Longleftrightarrow \text{non}(x\in A \text{ ou } x\in B).
\]
Par logique propositionnelle,
\[
\text{non}(P\text{ ou }Q)\Longleftrightarrow (\text{non }P)\text{ et }(\text{non }Q).
\]
Donc
\[
x\notin A\cup B
\Longleftrightarrow x\notin A \text{ et } x\notin B
\Longleftrightarrow x\in E\setminus A \text{ et } x\in E\setminus B.
\]
Ainsi
\[
x\in E\setminus(A\cup B)\Longleftrightarrow x\in (E\setminus A)\cap(E\setminus B),
\]
d’où l’égalité.
\end{proof}

\section{Produit cartésien}

\begin{definition}
Soient \(A\) et \(B\) deux ensembles. Le \textbf{produit cartésien} de \(A\) et \(B\) est l’ensemble
\[
A\times B=\{(a,b)\mid a\in A,\ b\in B\}.
\]
\end{definition}

\begin{exemple}
Si
\[
A=\{1,2\},
\qquad
B=\{x,y,z\},
\]
alors
\[
A\times B=\{(1,x),(1,y),(1,z),(2,x),(2,y),(2,z)\}.
\]
\end{exemple}

\begin{remarque}
L’ordre compte :
\[
(a,b)\neq (b,a)
\]
en général. Donc
\[
A\times B \neq B\times A
\]
comme ensembles d’objets ordonnés, même s’ils ont parfois le même nombre d’éléments lorsque \(A\) et \(B\) sont finis.
\end{remarque}

\section{Quantificateurs}

\subsection{Quantificateur universel}

\begin{definition}
L’écriture
\[
\forall x\in E,\ P(x)
\]
signifie : \textbf{pour tout} élément \(x\) de \(E\), la propriété \(P(x)\) est vraie.
\end{definition}

\begin{exemple}
\[
\forall x\in \R,\ x^2\geq 0.
\]
Cette phrase affirme que tout carré d’un réel est positif ou nul.
\end{exemple}

\subsection{Quantificateur existentiel}

\begin{definition}
L’écriture
\[
\exists x\in E,\ P(x)
\]
signifie : \textbf{il existe au moins un} élément \(x\) de \(E\) tel que \(P(x)\) soit vraie.
\end{definition}

\begin{exemple}
\[
\exists x\in\R,\ x^2=2.
\]
C’est vrai, par exemple pour \(x=\sqrt{2}\).
\end{exemple}

\subsection{Existence et unicité}

\begin{definition}
L’écriture
\[
\exists !\,x\in E,\ P(x)
\]
signifie : \textbf{il existe un unique} élément \(x\in E\) tel que \(P(x)\) soit vraie.
\end{definition}

\begin{exemple}
\[
\exists !\,x\in\R,\ 2x+3=7.
\]
En effet,
\[
2x=4 \Longrightarrow x=2.
\]
Il existe une solution, et elle est unique.
\end{exemple}

\begin{methode}
Pour prouver une assertion du type « il existe un unique \(x\) tel que \(P(x)\) », on sépare presque toujours la preuve en deux étapes :
\begin{enumerate}
    \item \textbf{Existence} : on exhibe un objet qui convient ;
    \item \textbf{Unicité} : on suppose que deux objets conviennent et on montre qu’ils sont égaux.
\end{enumerate}
\end{methode}

\begin{exemple}
Montrons que
\[
\exists !\,x\in\R,\ x+5=0.
\]

\textbf{Existence :} \(x=-5\) convient car
\[
-5+5=0.
\]

\textbf{Unicité :} si \(x\in\R\) vérifie \(x+5=0\), alors
\[
x=-5.
\]
Ainsi toute solution est égale à \(-5\), donc la solution est unique.
\end{exemple}

\subsection{Négation des quantificateurs}

\begin{proposition}
On a les équivalences logiques suivantes :
\[
\neg\big(\forall x\in E,\ P(x)\big)
\Longleftrightarrow
\exists x\in E,\ \neg P(x),
\]
\[
\neg\big(\exists x\in E,\ P(x)\big)
\Longleftrightarrow
\forall x\in E,\ \neg P(x).
\]
\end{proposition}

\begin{proof}
Dire que « il n’est pas vrai que tout \(x\) vérifie \(P(x)\) » signifie précisément qu’il existe au moins un élément pour lequel \(P(x)\) est fausse.

De même, dire que « il n’existe aucun \(x\) vérifiant \(P(x)\) » signifie que pour tout \(x\), la propriété \(P(x)\) est fausse.
\end{proof}

\begin{exemple}
La négation de
\[
\forall x\in\R,\ x^2+1>0
\]
est
\[
\exists x\in\R,\ x^2+1\leq 0.
\]

La négation de
\[
\exists x\in\R,\ x^2=-1
\]
est
\[
\forall x\in\R,\ x^2\neq -1.
\]
\end{exemple}

\section{Implication, réciproque, contraposée, équivalence}

\subsection{Implication}

\begin{definition}
Une implication est une proposition de la forme
\[
P \Longrightarrow Q,
\]
qui se lit « si \(P\), alors \(Q\) ».
\end{definition}

\begin{exemple}
\[
x=2 \Longrightarrow x^2=4.
\]
Cette implication est vraie.

En revanche,
\[
x^2=4 \Longrightarrow x=2
\]
est fausse, car on peut avoir \(x=-2\).
\end{exemple}

\subsection{Réciproque}

\begin{definition}
La \textbf{réciproque} de l’implication
\[
P\Longrightarrow Q
\]
est l’implication
\[
Q\Longrightarrow P.
\]
\end{definition}

\begin{remarque}
Une implication peut être vraie sans que sa réciproque le soit.
\end{remarque}

\subsection{Contraposée}

\begin{definition}
La \textbf{contraposée} de
\[
P\Longrightarrow Q
\]
est
\[
\neg Q \Longrightarrow \neg P.
\]
\end{definition}

\begin{proposition}
Une implication et sa contraposée sont logiquement équivalentes :
\[
P\Longrightarrow Q
\quad \Longleftrightarrow \quad
\neg Q \Longrightarrow \neg P.
\]
\end{proposition}

\begin{proof}
Supposons \(P\Longrightarrow Q\). Si \(\neg Q\) est vraie, alors \(P\) ne peut pas être vraie, sinon \(Q\) serait vraie. Donc \(\neg P\) est vraie. On a donc
\[
\neg Q \Longrightarrow \neg P.
\]

Réciproquement, supposons
\[
\neg Q \Longrightarrow \neg P.
\]
Montrons \(P\Longrightarrow Q\). Supposons \(P\) vraie. Si \(Q\) était fausse, alors \(\neg Q\) serait vraie, donc \(\neg P\) serait vraie, contradiction. Donc \(Q\) est vraie. Ainsi
\[
P\Longrightarrow Q.
\]
\end{proof}

\begin{methode}
Quand une implication semble difficile à démontrer directement, il faut penser à démontrer sa \textbf{contraposée}.
\end{methode}

\begin{exemple}
Montrons :
\[
x^2 \text{ impair } \Longrightarrow x \text{ impair}
\qquad (x\in \Z).
\]

Il est plus simple de raisonner par contraposée.

La contraposée est :
\[
x \text{ pair } \Longrightarrow x^2 \text{ pair}.
\]

Or si \(x\) est pair, il existe \(k\in\Z\) tel que \(x=2k\). Alors
\[
x^2=(2k)^2=4k^2=2(2k^2),
\]
donc \(x^2\) est pair. La contraposée est vraie, donc l’implication initiale aussi.
\end{exemple}

\subsection{Équivalence}

\begin{definition}
On dit que \(P\) et \(Q\) sont équivalentes si
\[
P\Longleftrightarrow Q,
\]
c’est-à-dire si l’on a à la fois
\[
P\Longrightarrow Q
\quad \text{et} \quad
Q\Longrightarrow P.
\]
\end{definition}

\begin{exemple}
Pour \(x\in\R\),
\[
x^2=0 \Longleftrightarrow x=0.
\]

En effet :
\begin{itemize}
    \item si \(x=0\), alors \(x^2=0\) ;
    \item si \(x^2=0\), alors nécessairement \(x=0\).
\end{itemize}
\end{exemple}

\section{Méthodes classiques de raisonnement}

\subsection{Raisonnement direct}

\begin{definition}
Le \textbf{raisonnement direct} consiste à partir des hypothèses pour en déduire la conclusion par une suite d’implications.
\end{definition}

\begin{exemple}
Montrons que la somme de deux entiers pairs est paire.

Soient \(a,b\in\Z\) deux entiers pairs. Il existe donc \(k,\ell\in\Z\) tels que
\[
a=2k,\qquad b=2\ell.
\]
Alors
\[
a+b=2k+2\ell=2(k+\ell).
\]
Comme \(k+\ell\in\Z\), on conclut que \(a+b\) est pair.
\end{exemple}

\subsection{Raisonnement par contraposée}

Nous l’avons déjà vu : pour montrer
\[
P\Longrightarrow Q,
\]
on peut montrer
\[
\neg Q \Longrightarrow \neg P.
\]

\begin{exemple}
Montrons :
\[
x^2\neq 0 \Longrightarrow x\neq 0.
\]

Cette implication se voit aussi directement, mais la contraposée est immédiate :
\[
x=0 \Longrightarrow x^2=0.
\]
Donc l’implication initiale est vraie.
\end{exemple}

\subsection{Raisonnement par l’absurde}

\begin{definition}
Le \textbf{raisonnement par l’absurde} consiste à supposer fausse la conclusion que l’on veut établir, puis à en déduire une contradiction.
\end{definition}

\begin{exemple}[Irrationalité de \(\sqrt{2}\)]
Montrons que \(\sqrt{2}\notin \Q\).

Supposons par l’absurde que \(\sqrt{2}\in \Q\). On peut alors écrire
\[
\sqrt{2}=\frac{p}{q}
\]
avec \(p\in\Z\), \(q\in\Z\setminus\{0\}\), et la fraction \(\frac{p}{q}\) irréductible.

En élevant au carré :
\[
2=\frac{p^2}{q^2}
\quad \Longrightarrow \quad
p^2=2q^2.
\]
Ainsi \(p^2\) est pair, donc \(p\) est pair. Il existe donc \(k\in\Z\) tel que
\[
p=2k.
\]
En remplaçant :
\[
(2k)^2=2q^2
\quad \Longrightarrow \quad
4k^2=2q^2
\quad \Longrightarrow \quad
q^2=2k^2.
\]
Donc \(q^2\) est pair, donc \(q\) est pair.

Ainsi \(p\) et \(q\) sont tous les deux pairs, donc ils ont \(2\) comme diviseur commun, ce qui contredit le fait que \(\frac{p}{q}\) soit irréductible.

La supposition de départ est absurde. Donc
\[
\sqrt{2}\notin \Q.
\]
\end{exemple}

\subsection{Raisonnement par disjonction de cas}

\begin{definition}
Le \textbf{raisonnement par cas} consiste à découper la situation en plusieurs possibilités exclusives, puis à traiter chacune d’elles.
\end{definition}

\begin{exemple}
Montrons que pour tout entier \(n\), le produit \(n(n+1)\) est pair.

Soit \(n\in\Z\). Deux cas sont possibles :
\begin{itemize}
    \item soit \(n\) est pair ; alors \(n(n+1)\) est pair ;
    \item soit \(n\) est impair ; alors \(n+1\) est pair, donc \(n(n+1)\) est pair.
\end{itemize}

Dans tous les cas, \(n(n+1)\) est pair.
\end{exemple}

\subsection{Raisonnement par récurrence}

\begin{theoreme}[Principe de récurrence]
Soit \(P(n)\) une propriété dépendant d’un entier naturel \(n\). Supposons que :
\begin{enumerate}
    \item \(P(0)\) est vraie ;
    \item pour tout \(n\in\N\), si \(P(n)\) est vraie alors \(P(n+1)\) est vraie.
\end{enumerate}
Alors \(P(n)\) est vraie pour tout \(n\in\N\).
\end{theoreme}

\begin{idee}
Le principe de récurrence fonctionne comme une rangée infinie de dominos :
\begin{itemize}
    \item on renverse le premier domino (initialisation) ;
    \item chaque domino renverse le suivant (hérédité) ;
    \item donc tous les dominos tombent.
\end{itemize}
\end{idee}

\begin{proof}
Admettons qu’il existe un entier \(n\in\N\) tel que \(P(n)\) soit fausse. Alors l’ensemble
\[
A=\{n\in\N\mid P(n)\text{ est fausse}\}
\]
est non vide. Comme \(\N\) est bien ordonné, \(A\) admet un plus petit élément \(n_0\).

On ne peut pas avoir \(n_0=0\), car \(P(0)\) est vraie.

Donc \(n_0\geq 1\), et ainsi \(n_0-1\in\N\). Par minimalité de \(n_0\), la propriété \(P(n_0-1)\) est vraie. Par hérédité, \(P(n_0)\) est alors vraie, contradiction.

Donc l’ensemble \(A\) est vide, ce qui signifie que \(P(n)\) est vraie pour tout \(n\in\N\).
\end{proof}

\begin{exemple}
Montrons que pour tout \(n\in\N\),
\[
1+2+\cdots+n=\frac{n(n+1)}{2}.
\]

\textbf{Initialisation :} pour \(n=0\),
\[
1+\cdots+0=0
\quad \text{et} \quad
\frac{0(0+1)}{2}=0.
\]
La propriété est vraie au rang \(0\).

\textbf{Hérédité :} supposons qu’au rang \(n\),
\[
1+2+\cdots+n=\frac{n(n+1)}{2}.
\]
Alors
\[
1+2+\cdots+n+(n+1)
=
\frac{n(n+1)}{2}+(n+1).
\]
On factorise par \(n+1\) :
\[
\frac{n(n+1)}{2}+(n+1)
=
(n+1)\left(\frac{n}{2}+1\right)
=
(n+1)\left(\frac{n+2}{2}\right)
=
\frac{(n+1)(n+2)}{2}.
\]
La propriété est donc vraie au rang \(n+1\).

Par récurrence, elle est vraie pour tout \(n\in\N\).
\end{exemple}
\subsection{Méthode d’analyse-synthèse}

\begin{definition}
La \textbf{méthode d’analyse-synthèse} est une méthode de recherche puis de démonstration, particulièrement utile dans les problèmes où l’on cherche à déterminer \textbf{tous} les objets vérifiant une certaine propriété.

\begin{remarque}
La méthode d’analyse-synthèse est particulièrement adaptée aux problèmes où l’on cherche à caractériser un objet par une propriété, c’est-à-dire à obtenir une condition nécessaire puis à montrer qu’elle est aussi suffisante.
\end{remarque}

Elle se décompose en deux étapes :
\begin{itemize}
    \item \textbf{Analyse :} on suppose que l’objet cherché existe, et on cherche quelles propriétés il doit nécessairement vérifier ;
    \item \textbf{Synthèse :} on repart des conditions obtenues, on construit effectivement un ou plusieurs objets, puis on vérifie qu’ils conviennent bien.
\end{itemize}
\end{definition}

\begin{remarque}
L’analyse est souvent une phase de \textbf{recherche} : elle permet de deviner la forme de la solution.

La synthèse est la phase de \textbf{rédaction rigoureuse} : elle permet de prouver que les objets trouvés conviennent effectivement.
\end{remarque}

\begin{idee}
On peut résumer la méthode ainsi :

\begin{center}
\textbf{On suppose que la solution existe} \(\longrightarrow\) \textbf{on découvre sa forme} \(\longrightarrow\) \textbf{on vérifie que cette forme convient réellement.}
\end{center}
\end{idee}

\begin{remarque}
Cette méthode est très fréquente en algèbre linéaire. On la retrouve par exemple :
\begin{itemize}
    \item pour déterminer toutes les solutions d’un système ;
    \item pour chercher l’écriture d’un vecteur dans une base ;
    \item pour trouver les applications linéaires vérifiant certaines conditions ;
    \item pour établir des résultats d’existence et d’unicité.
\end{itemize}
\end{remarque}

\begin{methode}
Lorsqu’un énoncé demande de \og déterminer tous les objets vérifiant une propriété \fg, on peut souvent suivre le schéma suivant :

\textbf{Analyse}
\begin{enumerate}
    \item On suppose qu’un objet \(x\) vérifie la propriété demandée ;
    \item On transforme les conditions pour obtenir des contraintes nécessaires sur \(x\) ;
    \item On devine alors une forme possible des solutions.
\end{enumerate}

\textbf{Synthèse}
\begin{enumerate}
    \item On considère les objets obtenus pendant l’analyse ;
    \item On vérifie qu’ils satisfont bien toutes les conditions ;
    \item On conclut que ce sont exactement les solutions cherchées.
\end{enumerate}
\end{methode}

\begin{exemple}[Résolution d’une équation par analyse-synthèse]
Déterminons tous les réels \(x\) tels que
\[
x+\frac{1}{x} = 2.
\]

\textbf{Analyse :} supposons qu’il existe \(x\in\R\) vérifiant cette équation. Comme \(\frac{1}{x}\) apparaît, on doit avoir
\[
x\neq 0.
\]
On multiplie alors l’équation par \(x\) :
\[
x^2+1=2x.
\]
On réorganise :
\[
x^2-2x+1=0.
\]
On reconnaît une identité remarquable :
\[
(x-1)^2=0.
\]
Donc nécessairement,
\[
x=1.
\]

Ainsi, s’il existe une solution, elle ne peut être que \(x=1\).

\textbf{Synthèse :} vérifions que \(x=1\) convient :
\[
1+\frac{1}{1}=1+1=2.
\]
Donc \(x=1\) est bien solution.

On conclut que l’ensemble des solutions est
\[
\{1\}.
\]
\end{exemple}

\begin{exemple}[Exemple type d’existence et unicité]
Montrons que pour tout réel \(a\), il existe un unique réel \(x\) tel que
\[
2x+a=0.
\]

\textbf{Analyse :} supposons qu’un réel \(x\) vérifie
\[
2x+a=0.
\]
Alors nécessairement
\[
2x=-a,
\qquad
x=-\frac{a}{2}.
\]
Ainsi, s’il existe une solution, elle ne peut être que
\[
x=-\frac{a}{2}.
\]

\textbf{Synthèse :} vérifions que ce réel convient effectivement :
\[
2\left(-\frac{a}{2}\right)+a=-a+a=0.
\]
Donc \(x=-\frac{a}{2}\) est bien solution.

Enfin, comme toute solution éventuelle est nécessairement égale à \(-\frac{a}{2}\), cette solution est unique.

On a donc montré :
\[
\exists !\,x\in\R,\ 2x+a=0.
\]
\end{exemple}

\begin{remarque}
Dans beaucoup de problèmes, l’analyse ne constitue pas à elle seule une preuve complète.  
Elle permet seulement de trouver ce que la solution \textbf{doit être}. La preuve n’est terminée qu’après la synthèse, c’est-à-dire après la vérification explicite.
\end{remarque}

\begin{exemple}[Exemple très simple en géométrie ou en algèbre linéaire]
Si l’on cherche un vecteur \(u\) tel que
\[
u+v=w,
\]
on peut raisonner par analyse-synthèse :

\textbf{Analyse :} si un tel vecteur \(u\) existe, alors nécessairement
\[
u=w-v.
\]

\textbf{Synthèse :} on vérifie que le vecteur \(u=w-v\) convient bien :
\[
(w-v)+v=w.
\]

Donc il existe un unique vecteur \(u\) vérifiant \(u+v=w\), à savoir
\[
u=w-v.
\]
\end{exemple}

\section{Existence et unicité : une méthode centrale}

\begin{remarque}
En algèbre linéaire, on rencontre en permanence des énoncés du type :
\begin{itemize}
    \item « tout vecteur s’écrit de manière unique comme combinaison linéaire de la base » ;
    \item « il existe une unique application linéaire vérifiant telle propriété » ;
    \item « tout système admet une unique solution dans telle situation ».
\end{itemize}

Il faut donc bien comprendre la structure d’une preuve d’existence et d’unicité.
\end{remarque}

\begin{methode}
Pour démontrer :
\[
\exists !\,x\in E,\ P(x),
\]
on fait :

\textbf{Étape 1 : existence}  
On construit explicitement un objet \(x\) et on vérifie qu’il satisfait \(P(x)\).

\textbf{Étape 2 : unicité}  
On suppose que \(x\) et \(y\) vérifient tous les deux \(P\), puis on montre que \(x=y\).
\end{methode}

\begin{exemple}
Montrons que pour tout réel \(a\), il existe un unique réel \(x\) tel que
\[
x+a=0.
\]

\textbf{Existence :} le réel \(x=-a\) convient car
\[
(-a)+a=0.
\]

\textbf{Unicité :} si \(x\) vérifie \(x+a=0\), alors
\[
x=-a.
\]
Toute solution est donc nécessairement égale à \(-a\). Il y a donc unicité.
\end{exemple}

\section{Erreurs classiques à éviter}

\begin{itemize}
    \item Confondre \(x\in A\) et \(x\subset A\).  
    Un élément appartient à un ensemble ; un ensemble est inclus dans un autre.
    
    \item Pour montrer \(A=B\), ne prouver qu’une seule inclusion.  
    Il faut presque toujours montrer les deux inclusions.

    \item Confondre implication et équivalence.  
    Le fait que \(P\Longrightarrow Q\) soit vraie ne signifie pas que \(Q\Longrightarrow P\) soit vraie.

    \item Mal nier une proposition quantifiée.  
    Par exemple, la négation de
    \[
    \forall x,\ P(x)
    \]
    n’est pas
    \[
    \forall x,\ \neg P(x),
    \]
    mais
    \[
    \exists x,\ \neg P(x).
    \]

    \item Oublier de vérifier l’unicité dans une question d’existence et unicité.

    \item En récurrence, oublier l’initialisation ou écrire une hérédité incomplète.
\end{itemize}

\section{Résumé du chapitre}

\begin{itemize}
    \item Un ensemble est une collection d’éléments.
    \item \(A\subset B\) signifie : tout élément de \(A\) appartient à \(B\).
    \item Pour montrer \(A=B\), on montre \(A\subset B\) puis \(B\subset A\).
    \item Les opérations fondamentales sont : union, intersection, différence, complémentaire.
    \item Les quantificateurs \(\forall\), \(\exists\), \(\exists !\) doivent être maniés avec rigueur.
    \item Une implication n’est pas une équivalence.
    \item La contraposée d’une implication lui est logiquement équivalente.
    \item Les méthodes fondamentales de preuve sont :
    \begin{itemize}
        \item raisonnement direct,
        \item contraposée,
        \item absurde,
        \item disjonction de cas,
        \item récurrence.
    \end{itemize}
    \item Les preuves d’existence et unicité jouent un rôle central dans tout le reste du cours.
\end{itemize}

\section*{Exercices d’application immédiate}

\begin{enumerate}[label=\textbf{Exercice \arabic*.}]
    \item Montrer que pour tous ensembles \(A,B,C\),
    \[
    A\cap(B\setminus C)=(A\cap B)\setminus C.
    \]

    \item Nier correctement les propositions suivantes :
    \begin{enumerate}
        \item \(\forall x\in\R,\ x^2+x+1>0\),
        \item \(\exists x\in\R,\ e^x=0\),
        \item \(\forall n\in\N,\ \exists p\in\N,\ p\geq n\).
    \end{enumerate}

    \item Montrer par récurrence que pour tout \(n\in\N\),
    \[
    1+3+5+\cdots +(2n+1)=(n+1)^2.
    \]

    \item Montrer que si \(n^2\) est multiple de \(3\), alors \(n\) est multiple de \(3\).

    \item Montrer que pour tout réel \(a\), l’équation
    \[
    x-a=0
    \]
    admet une unique solution.
\end{enumerate}

\end{document}